\noindent
\textbf{One-line summary.} Generative neural memory enabling fast and efficient use of past experience. We demonstrate its effectiveness for in-context RL.\\
\textbf{Real-world demo.} Adaptive humanoid (physical system of adaptive agent paper~\cite{team2023human}), see example in paragraph 1 below.

\section{Overview}
Long-horizon autonomy requires robots to retain and recall prior experience and reuse acquired knowledge. For example, a robot trying to open a high cabinet first fails to reach it, even after jumping, then remembers these failures, infers the cabinet is out of reach, and later uses a box as a step to reach and open it. To effectively push the box, it may also reuse prior interaction experience to estimate the force needed to move it. We refer to such storage, recall, and use of past experience over extended periods of time ($\geq$ 10 minutes) as long-term memory.

Prior work adds long-term memory to transformers via KV stores~\cite{wu2022memorizing,borgeaud2022improving, vennerod2021long} with search-based retrieval; however, retrieval latency scales with memory size and time, hindering real-time robotic control. A separate line of work stores history in a fixed-size recurrent state~\cite{dai2019transformer,bulatov2022recurrent,rodkin2024associative,munkhdalai2024leave,xiao2023efficient,gu2024mamba,vennerod2021long}, which limits information capacity due to the state dimension.
We instead propose to augment transformers with learned neural memory~\cite{behrouz2024titans,sun2024learning}, using the network itself as a compressor to store much more while keeping per-step retrieval constant.
However, most neural memory approaches couple memory with a predictive readout that yields a single point estimate, limiting counterfactual reasoning. We instead propose a generative neural memory, akin to humans~\cite{spens2024generative,bennett2023brief}, that can sample diverse variations of relevant experiences, enabling faster learning during long-horizon control (see below for differences between a world model and neural memory).

Our goal is for the robot to exploit past experience to solve the current task; the experience may come from earlier in the same episode or from prior episodes.

\textbf{Test-time.}
At each step, every Transformer block forms a query from its current features and consults its \emph{per-block generative neural memory}. The memory returns one or more latent variants that are integrated by cross-attention to produce memory-augmented block outputs, which the policy uses to act. After the transition, the fast weights are updated by a small write step, so subsequent queries reflect the latest information. This loop---query, sample, attend, act, write---runs throughout the episode.

\textbf{Training-time.}
We meta-train across tasks with the same read-write loop active inside each rollout. The base weights of the policy and memory are optimized end-to-end to maximize return while regularizing the memory's posterior against its prior; the fast weights adapt online but are reset between tasks.
The \emph{outer loop} teaches the model \emph{how to} structure queries, priors, and writes, while the \emph{inner loop}, small fast-weight updates, \emph{does} the rapid task-specific adaptation.

\begin{tcolorbox}[clarifications]
\textbf{Clarification 1: Distinction between world-model and generative neural memory}
\\\\
World model is an explicit dynamics function \(f_{\theta}\) mapping a state–action pair to the next state, e.g., \(s_{t+1} = f_{\theta}(s_t, a_t)\).\\\\
Neural memory is a differentiable key-value store updated online, possibly via gradient-based writes. The values stored are not limited to dynamics: they may be latent vectors \(v_t = W_v x_t\) or raw transitions such as \((s_{t-1}, a_t, s_t)\). This makes neural memory strictly more general than a dynamics-only model: it can emulate a model by caching transitions and, additionally, may hold task context, affordances, goals, or other non-dynamics information.
\end{tcolorbox}

\begin{tcolorbox}[clarifications]
\textbf{Clarification 2: Why in-context RL?}
\\\\
Our objective is to demonstrate the efficiency of a generative neural memory. In-context RL is a natural testbed: the agent must explore, gather information, and adapt its behavior, directly evaluating our hypothesis that such a memory enables rapid, sample-efficient learning.
\end{tcolorbox}

\begin{tcolorbox}[clarifications]
\textbf{Clarification 3: Why is the real-world system a humanoid?}
\\\\
The sim-to-real and RL ecosystem for humanoids is particularly mature, enabling us to train an in-context RL policy entirely in simulation and deploy it on hardware. Prior work using motion-tracking--based reward shaping shows that tasks of comparable complexity can transfer successfully to the real world.
\\\\
The extension is non-trivial with RL on a high-dimensional action space. See Experimental Plan for details.
\end{tcolorbox}

\begin{tcolorbox}[clarifications]
\textbf{Clarification 4: Is this one or two different papers?}
\\\\
Ideally, we would want this one paper to make it a very strong submission. The reward designs for a humanoid can be time-consuming (and sim2real). We have three months to execute before RSS.
\end{tcolorbox}